\clearpage\thispagestyle{empty}
\pdfbookmark[0]{Quellen zur Besprechung der Connoissance des tems für 1812}{connoissance1812-sources}
\begin{center}\textbf{Quellen zur Besprechung der Connoissance des tems für 1812}\end{center}
\noindent Diese KI-Übersetzung ins heutige Deutsch folgt der editierbaren Transkription von Carl Friedrich Gauß, \textit{Werke}, Bd.~VI (Göttingen, 1874), S.~537--541. Die Besprechung erschien zuerst unter Paris in den \textit{Göttingischen gelehrten Anzeigen}, Stück~199, 14. Dezember 1811, S.~1978--1987. Das Institut de mécanique céleste et de calcul des éphémérides verzeichnet den besprochenen \href{https://cdt.imcce.fr/items/show/161}{Jahrgang für 1812}.
\par\smallskip\noindent Die Jahreszahl nach \textit{Juillet} erscheint im Druck der \textit{Werke} als 181c und im digitalisierten Erstdruck als 1310. Die Übersetzung behält 181c bei. Im Scan der \textit{Werke} steht beim ersten Delambre-Beitrag \textit{von Gauss}. Im Titel der \textit{Theoria motus} scheint dort \textit{carporum} zu stehen. Die editierbare Transkription hat \textit{vom Gauss} und \textit{corporum}.
\par\smallskip\noindent Die Besprechung nennt das Septemberheft 1808 der \textit{Monatlichen Correspondenz}. Der zugehörige Aufsatz in den \textit{Werken}, Bd.~VI, S.~129, nennt im Kopf Oktober 1808. Das letzte Glied des Ausdrucks für den Fehler der Polhöhe auf S.~139 des Aufsatzes hat ein Minuszeichen. In der Besprechung steht ein Pluszeichen.
\par\smallskip\noindent Das erste französische Zitat besagt, Burckhardt habe \emph{alle} Beobachtungen, die alten wie die neueren, verglichen und neu berechnet. Laut Delambres erster Bemerkung beweise die letzte Formel, dass die Methode in der Praxis nur eine recht begrenzte Genauigkeit hätte. Nach seiner zweiten Bemerkung würde man in einer Nacht, in der drei Sterne in gleicher Höhe beobachtet werden könnten, sehr wahrscheinlich einen bekannten Stern durch den Meridian gehen sehen und so mit weniger Mühe und größerer Genauigkeit die Breite und aus einer einzigen Höhe die Zeit der Pendeluhr erhalten.
\par\smallskip\noindent \href{https://gdz.sub.uni-goettingen.de/id/PPN236007467}{Text der Werkausgabe}. Erstdruck, \href{https://images.p.sub.uni-goettingen.de/iiif/image/gdz:PPN319721507_1811:00002012/full/1650,/0/default.jpg}{S.~1978}, \href{https://images.p.sub.uni-goettingen.de/iiif/image/gdz:PPN319721507_1811:00002013/full/1650,/0/default.jpg}{S.~1979}, \href{https://images.p.sub.uni-goettingen.de/iiif/image/gdz:PPN319721507_1811:00002014/full/1650,/0/default.jpg}{S.~1980}, \href{https://images.p.sub.uni-goettingen.de/iiif/image/gdz:PPN319721507_1811:00002015/full/1650,/0/default.jpg}{S.~1981}, \href{https://images.p.sub.uni-goettingen.de/iiif/image/gdz:PPN319721507_1811:00002016/full/1650,/0/default.jpg}{S.~1982}, \href{https://images.p.sub.uni-goettingen.de/iiif/image/gdz:PPN319721507_1811:00002017/full/1650,/0/default.jpg}{S.~1983}, \href{https://images.p.sub.uni-goettingen.de/iiif/image/gdz:PPN319721507_1811:00002018/full/1650,/0/default.jpg}{S.~1984}, \href{https://images.p.sub.uni-goettingen.de/iiif/image/gdz:PPN319721507_1811:00002019/full/1650,/0/default.jpg}{S.~1985}, \href{https://images.p.sub.uni-goettingen.de/iiif/image/gdz:PPN319721507_1811:00002020/full/1650,/0/default.jpg}{S.~1986} und \href{https://images.p.sub.uni-goettingen.de/iiif/image/gdz:PPN319721507_1811:00002021/full/1650,/0/default.jpg}{S.~1987}.
\par\smallskip\noindent Übersetzung und Satz wurden von KI erstellt.
\clearpage
